\begin{table}[htbp]
\centering
\caption[Momentos estadísticos y correlación de PC1--PC3 de precipitación - Subconjunto p90]
{Asimetría ($\gamma_1$) y curtosis ($\gamma_2$) de la PC1 de precipitación, y correlación de Spearman ($r$) entre PC1, PC2 y PC3 (con signo) y la vorticidad relativa a 850 hPa asociada a la fase (con signo, negativa en el Hemisferio Sur), para el subconjunto extremo (p90). El signo de la correlación depende de la convención de signo de cada EOF (Sección~\ref{subsec:analisis_pc}); lo relevante es que el lado del patrón dominante en que cae cada ciclón está asociado a su intensidad. Los valores significativos ($p<0.05$) se muestran en negrita.}
\label{tab:stat_p90_tp}
\begin{tabular*}{\textwidth}{@{\extracolsep{\fill}}llrrrrr}
\hline
\textbf{Región} & \textbf{Fase} & \textbf{Asimetría ($\gamma_1$)} & \textbf{Curtosis ($\gamma_2$)} & \textbf{PC1} & \textbf{PC2} & \textbf{PC3} \\
\hline
\textbf{SBR} & Incipiente       & 1.57 & 3.48  & \textbf{-0.36} & \textbf{-0.53} & -0.22 \\
             & Intensificación & 1.20 & 1.67  & \textbf{-0.48} & \textbf{-0.43} & -0.24 \\
             & Madurez         & 0.41 & -0.67 & -0.17 & -0.22 & -0.11 \\
             & Decaimiento     & 4.02 & 19.77 & \textbf{-0.34} & 0.17 & -0.18 \\
\hline
\textbf{LPB} & Incipiente       & 1.90 & 4.52  & \textbf{-0.49} & \textbf{-0.31} & 0.22 \\
             & Intensificación & 0.60 & -0.38 & \textbf{-0.43} & -0.01 & \textbf{0.26} \\
             & Madurez         & 0.56 & 0.16  & -0.03 & -0.18 & 0.02 \\
             & Decaimiento     & 4.34 & 23.18 & \textbf{-0.50} & 0.09 & 0.03 \\
\hline
\textbf{ARG} & Incipiente       & 4.43 & 30.88 & \textbf{-0.44} & \textbf{-0.29} & \textbf{-0.23} \\
             & Intensificación & 1.87 & 5.15  & \textbf{-0.37} & \textbf{0.21} & \textbf{-0.28} \\
             & Madurez         & 1.66 & 4.53  & \textbf{-0.24} & 0.03 & -0.14 \\
             & Decaimiento     & 1.96 & 3.37  & \textbf{-0.39} & \textbf{0.32} & 0.10 \\
\hline
\end{tabular*}
\end{table}
